% NOTE: Filename is fixed by the book outline; content teaches Week 16
% (Business Intelligence with Amazon QuickSight & Reverse ETL), plus the
% Milestone 4 capstone.
\chapter{Business Intelligence with QuickSight and Reverse ETL}
\index{Amazon QuickSight}\index{SPICE}\index{row-level security}\index{reverse ETL}\index{UNLOAD}\index{identity resolution}\index{anomaly detection}%
\begin{objectives}
\begin{itemize}
    \item Explain Amazon QuickSight architecture and the SPICE in-memory engine's role and limits.
    \item Connect QuickSight to Redshift, Athena, and S3 with appropriate refresh and direct-query trade-offs.
    \item Enforce row-level security on shared datasets and apply ML Insights for anomaly detection.
    \item Implement reverse ETL: Redshift UNLOAD to S3, Lambda push to a CRM/marketing API, with identity resolution upstream.
    \item Reason about sync-frequency trade-offs: freshness versus API rate limits and cost.
    \item Design the BI access layer for 500 analysts with region-scoped visibility, and assemble the Milestone~4 real-time fraud capstone.
\end{itemize}
\end{objectives}

\begin{crosscloud}
The serving layer---dashboards for humans, operational pushes back to business systems---closes the analytics loop on every platform.

\vspace{2mm}
{\footnotesize
\begin{tabular}{@{}p{2.8cm}p{3.0cm}p{3.0cm}p{3.0cm}@{}}
\toprule
\textbf{Capability} & \textbf{AWS} & \textbf{Azure} & \textbf{GCP} \\
\midrule
Managed BI & Amazon QuickSight & Power BI (Fabric) & Looker / Looker Studio \\
In-memory acceleration & SPICE & Power BI import/semantic models & BI Engine \\
Row-level security & QuickSight RLS rules & Power BI RLS & Looker access filters \\
ML-assisted insights & QuickSight ML Insights & Copilot / Smart visuals & Gemini in Looker \\
Warehouse export & Redshift UNLOAD & COPY INTO / Fabric export & BigQuery EXPORT DATA \\
Reverse-ETL glue & Lambda + API calls & Logic Apps / Functions & Cloud Run functions \\
\addlinespace
Reader pricing & Per-reader session & Per-user / capacity & Per-user \\
Sharing unit & Dashboard / analysis & Workspace app / semantic model & Look / dashboard \\
\bottomrule
\end{tabular}}

\vspace{1mm}
\textbf{Snowflake} hosts dashboards in Snowsight and streams reverse ETL through external functions; the ecosystem (Census, Hightouch) serves all platforms. \textbf{MongoDB} Charts covers operational BI directly on collections.
\end{crosscloud}

\section{Week 16 Roadmap and Mental Model}
Everything before this week moved data toward analytical truth; Week~16 moves it toward \emph{decisions and actions}. That happens in two directions. Outward to humans: dashboards, reports, and governed self-service---Amazon QuickSight's territory. And outward to systems: pushing curated segments and scores back into the operational tools where work happens---reverse ETL.

The engineering shift is from correctness to \emph{governance at scale}. A dashboard seen by 500 analysts with different entitlements is an access-control system with charts attached. A reverse-ETL pipeline writing to a CRM is a production integration with rate limits, identity ambiguity, and business consequences. This week treats both with the same rigor as the ingestion side \cite{awsQuickSightDocs}.

\begin{definitionbox}
\textbf{Reverse ETL} is the practice of syncing curated warehouse data---segments, scores, aggregates---back into operational systems (CRMs, marketing platforms, support tools) so that analytics directly drives business processes.
\end{definitionbox}

\begin{keyterms}
\textbf{SPICE}: QuickSight's in-memory calculation engine; datasets imported into SPICE serve fast, cheap dashboard queries. \textbf{Direct query}: QuickSight querying the source (Redshift/Athena) live instead of SPICE. \textbf{RLS}: row-level security via a rules dataset mapping users/groups to allowed dimension values. \textbf{ML Insights}: QuickSight's built-in anomaly detection, forecasting, and auto-narratives. \textbf{UNLOAD}: Redshift's parallel export command to S3. \textbf{Identity resolution}: matching records for the same person across systems before syncing.
\end{keyterms}

\section{Monday: QuickSight Architecture and SPICE}
\subsection{The Serving Layer, Managed}
QuickSight is serverless BI: no servers, per-reader pricing, browser-based authoring, and embedding APIs for application dashboards. A \emph{dataset} defines a source and its preparation; an \emph{analysis} is the author-facing canvas; a \emph{dashboard} is the published, shared, read-only product. Because it is serverless, the capacity question moves from ``how many nodes'' to ``SPICE or direct query?''
\index{QuickSight}

\subsection{SPICE: The Acceleration Decision}
SPICE (Super-fast, Parallel, In-memory Calculation Engine) imports dataset snapshots into managed memory. Dashboards against SPICE respond in sub-seconds without touching the source, which is both a performance and a cost story: a thousand dashboard loads do not become a thousand Redshift queries. The limits are the dataset capacity per account and the \emph{refresh} cadence---SPICE data is only as fresh as its last refresh, so freshness-hungry dashboards either schedule tight refreshes (which do hit the source) or use direct query and pay the source for every interaction \cite{awsQuickSightDocs}.
\index{SPICE}

\begin{table}[H]
\centering
\small
\begin{tabular}{@{}p{3.4cm}p{4.4cm}p{4.4cm}@{}}
\toprule
\textbf{Criterion} & \textbf{SPICE import} & \textbf{Direct query} \\
\midrule
Freshness & As of last refresh & Live \\
Dashboard latency & Sub-second, load-independent & Source-dependent \\
Load on source & Refresh schedule only & Every interaction \\
Capacity limit & SPICE allocation per account & Source concurrency \\
Best for & Shared executive/operational dashboards & Live operational monitoring \\
\bottomrule
\end{tabular}
\caption{SPICE versus direct query.}
\label{tab:ch16-spice}
\end{table}

\begin{pitfall}
\textbf{Defaulting everything to direct query.} Five hundred analysts on live Redshift queries re-creates the Chapter~6 concurrency problem at the BI layer. Default to SPICE for shared dashboards; reserve direct query for the few views whose freshness requirement is measured in minutes.
\end{pitfall}

\section{Tuesday: Connecting QuickSight to Redshift, Athena, and S3}
QuickSight connects natively to the week's familiar sources. \textbf{Redshift} connections support both SPICE and direct query and can target Serverless workgroups. \textbf{Athena} connections query the Glue catalog directly---the governed lake from Chapters~7--8 becomes a dashboard source with no data movement, and Lake Formation column masks apply through the Athena path. \textbf{S3} manifests allow file-set datasets for simple cases. Cross-account patterns matter here: dashboards in a central BI account can query Redshift datashares (Chapter~8) and catalogs shared via RAM, keeping data in domain accounts while the BI layer centralizes.
\index{Athena data source}\index{cross-account BI}

\begin{tipbox}
Give QuickSight its own Redshift user or IAM-identity-mapped role and a dedicated WLM service class. Then the Chapter~6 telemetry answers ``what did the BI layer do to the warehouse at 9:05 a.m.?'' without archaeology.
\end{tipbox}

\section{Wednesday: Row-Level Security, ML Insights, and Reverse ETL}
\subsection{Row-Level Security on Shared Datasets}
QuickSight RLS attaches a \emph{rules dataset} to a shared dataset: rows map users or groups to allowed values of a column (\texttt{region = 'EU'} for the EU managers group). Every visual, filter, and export from that dataset applies the restriction. The alternative---duplicating datasets per region---is a maintenance failure waiting for its second region. RLS keeps one dataset, one definition of the metrics, and per-viewer data \cite{awsQuickSightDocs}.
\index{row-level security}

\subsection{ML Insights}
QuickSight's built-in ML features---anomaly detection on time series, forecasting, and auto-generated narratives---run on SPICE datasets without a modeling pipeline. For data engineers they are the fastest way to put a drift tripwire on a business metric: configure anomaly detection on daily revenue, and the first day the pipeline breaks, the \emph{business} metric alarms too, complementing the infrastructure alarms from earlier chapters.
\index{ML Insights}\index{anomaly detection}

\subsection{Reverse ETL: Closing the Loop}
Reverse ETL pushes warehouse truth back to operational tools. The AWS-native shape: \textbf{UNLOAD} a curated segment query from Redshift to S3 \cite{awsRedshiftUNLOAD}, an S3 event triggers a \textbf{Lambda} that batches records to the CRM/marketing API with rate-limit-aware throttling and dead-letter capture, and a run-log table records what was synced when.

Two disciplines separate a pipeline from a liability. \textbf{Identity resolution} comes first: the warehouse row must map to the right person in the CRM---matching on a stable identity key established upstream, not on email strings that differ in case and era. And \textbf{sync frequency} is a budget: hourly syncs against a rate-limited API can cost more in calls than the value of the freshness; state the freshness requirement, then engineer down to it.
\index{reverse ETL}\index{identity resolution}\index{sync frequency}

\begin{lstlisting}[language=SQL,caption={UNLOAD a churn-risk segment for reverse ETL.},label={lst:ch16-unload}]
UNLOAD ($$
  SELECT customer_id, churn_score, segment
  FROM analytics.customer_scores
  WHERE score_date = CURRENT_DATE AND churn_score > 0.7
$$)
TO 's3://acme-analytics-raw-123456789012/reverse-etl/churn/'
IAM_ROLE 'arn:aws:iam::123456789012:role/RedshiftS3ReadRole'
FORMAT AS CSV
HEADER
PARALLEL OFF;
\end{lstlisting}

\begin{pitfall}
\textbf{Syncing duplicates to systems of engagement.} A warehouse join bug that double-counts is embarrassing; the same bug reverse-ETLed emails a customer twice or double-books a discount. Treat reverse ETL as a production integration: idempotency keys on the push, dry-run modes, and reconciliation counts between UNLOAD output and API acknowledgments.
\end{pitfall}

\section{Thursday Hands-on Project: A Monitored Sales Dashboard}
\index{hands-on project!QuickSight dashboard}
Connect QuickSight to the Redshift warehouse from Part~2, build a sales monitoring dashboard, apply a filter, and configure an ML anomaly-detection alert on the revenue trend.

\subsection*{Lab Objectives}
\begin{itemize}
    \item Create a QuickSight dataset from Redshift over \texttt{fact\_sales} and its dimensions.
    \item Build a dashboard with a KPI, a trend line, and a working region filter.
    \item Configure anomaly detection on daily revenue with an alert to email/SNS.
    \item Compare SPICE import and direct query on the same dataset.
\end{itemize}

\subsection*{Procedure}
\begin{enumerate}
    \item In QuickSight, create a dataset from the Redshift Serverless workgroup: a join of \texttt{fact\_sales} to \texttt{dim\_date}, \texttt{dim\_product}, and \texttt{dim\_store}.
    \item Choose SPICE import; schedule a daily refresh after the ETL window.
    \item Build the analysis: KPI of month-to-date revenue, line chart of daily revenue, bar of revenue by category, and a \texttt{region} filter control affecting all visuals.
    \item Add an anomaly-detection insight on daily revenue and configure a threshold alert.
    \item Duplicate the dataset as direct query; compare dashboard interaction latency and observe the Redshift query log entries per interaction.
    \item Publish as a dashboard shared with a test reader group.
\end{enumerate}

\subsection*{Verification}
\begin{itemize}
    \item The region filter restricts every visual; exported data respects it.
    \item The anomaly alert fires when a synthetic revenue dip is injected into the fact table and the refresh runs.
    \item The SPICE dashboard is materially faster than direct query under ten concurrent viewers; the direct-query variant shows per-interaction queries in \texttt{SVL\_QUERY\_SUMMARY}.
    \item Refresh timestamps on the dashboard match the scheduled cadence.
\end{itemize}

\section{Friday Use Case: Governed BI for 500 Analysts}
\index{use case!enterprise BI access layer}
An enterprise rolls out QuickSight to 500 analysts across five regions. Every analyst queries centralized warehouse data; regional sales managers must see only their region's rows. Data engineering owns the datasets; analysts own their analyses. Design the access layer.

\subsection{Design Decisions}
\textbf{One governed dataset per domain.} Curated, documented datasets (sales, inventory, marketing) are built by data engineering on conformed warehouse marts; analysts build analyses on top but cannot redefine the metrics layer. This is the semantic-layer bargain: consistency in exchange for a backlog.

\textbf{RLS from the identity source.} The rules dataset maps IAM-Identity-Center groups (\texttt{sales-eu}, \texttt{sales-na}, \ldots) to region values. Group membership is managed in the identity provider, so access reviews and joiner/mover/leaver processes inherit the existing HR workflow---no BI-side access list to drift.

\textbf{SPICE with staggered refreshes.} All shared datasets use SPICE, refreshed after the nightly ETL and again midday for the operational domain. Direct query is an exception requiring a freshness justification, protecting the warehouse from the 9:05 a.m.\ refresh storm.

\textbf{Cost visibility.} Per-reader session pricing makes adoption cheap to start; the platform team dashboards monthly reader-session counts per department so finance sees BI cost as a function of actual usage.

\begin{table}[H]
\centering
\small
\begin{tabular}{@{}p{3.6cm}p{4.4cm}p{4.2cm}@{}}
\toprule
\textbf{Requirement} & \textbf{Mechanism} & \textbf{Failure it prevents} \\
\midrule
Region-scoped visibility & RLS rules dataset by IdP group & Per-region dataset forks \\
Consistent metric definitions & Governed datasets only & Competing revenue numbers \\
Warehouse protection & SPICE default, staggered refresh & 9:05 a.m.\ query storms \\
Access lifecycle & Group membership via IdP & Orphaned access after transfers \\
Adoption cost control & Per-reader session reporting & Silent license sprawl \\
\bottomrule
\end{tabular}
\caption{The 500-analyst BI access layer.}
\label{tab:ch16-bi-access}
\end{table}

\begin{pitfall}
\textbf{RLS rules managed by hand in a spreadsheet.} The rules dataset \emph{is} the authorization system for the BI layer. Derive it from the identity provider's groups with a small pipeline, review its diffs like code, and alert when a user appears in no rule (fail-closed: they see nothing) versus in a wildcard rule (fail-open: they see everything).
\end{pitfall}

\section{Interview Q\&A}
\subsection{Concepts and Trade-offs}
\begin{enumerate}
    \item \textbf{Q: What does SPICE accelerate, and what is its capacity limit?}\\
    A: Dataset queries for dashboards---imported in-memory snapshots answer in sub-seconds without touching the source. Capacity is bounded by the SPICE allocation per account, and freshness by the refresh schedule.
    \item \textbf{Q: How does QuickSight row-level security filter a shared dataset?}\\
    A: A rules dataset maps users/groups to allowed column values; every visual and export applies the restriction at query time, so one dataset serves many audiences.
    \item \textbf{Q: What is reverse ETL, and which operational systems does it feed?}\\
    A: Syncing curated warehouse results---segments, scores---back into operational tools: CRMs, marketing platforms, support systems. Analytics becomes action inside the systems where work happens.
    \item \textbf{Q: Why must identity resolution happen before syncing audiences to a CRM?}\\
    A: The warehouse key must map to the correct person in the target system; syncing on uncontrolled identifiers creates duplicate or mis-targeted outreach with real business consequences.
    \item \textbf{Q: Hourly versus daily sync: which trade-offs decide the frequency?}\\
    A: API rate limits and call cost against the business value of freshness. State the freshness requirement first; then engineer to the loosest cadence that meets it.
    \item \textbf{Q: When is direct query justified over SPICE?}\\
    A: When freshness requirements are measured in minutes and the source can bear per-interaction load---typically operational monitors with small audiences, not company-wide dashboards.
\end{enumerate}

\section{Further Reading}
The QuickSight User Guide covers SPICE, data sources, RLS rules datasets, ML Insights, and embedding \cite{awsQuickSightDocs}. The UNLOAD reference covers parallel export options for reverse ETL \cite{awsRedshiftUNLOAD}; the Lambda guide covers the throttled-push pattern to external APIs \cite{awsLambdaDocs}. The Well-Architected Analytics Lens frames the governed-self-service balance \cite{awsWellArchitectedAnalytics}.

\section*{Key Takeaways}
\begin{itemize}
    \item QuickSight makes BI serverless; the real capacity decision is SPICE versus direct query per dataset.
    \item Governed datasets plus RLS rules let one semantic layer serve every audience safely.
    \item ML Insights puts anomaly tripwires on business metrics---complementing, not replacing, infrastructure alarms.
    \item Reverse ETL turns the warehouse into an action system; identity resolution and sync budgets make it safe.
    \item Treat CRM pushes as production integrations: idempotency, dry runs, reconciliation.
    \item The BI layer's access controls derive from the identity provider, not from dashboard-side lists.
\end{itemize}

\section{Exercises}
\begin{enumerate}
    \item \textbf{(Conceptual)} Explain why SPICE both reduces warehouse load and creates a freshness contract with the business.
    \item \textbf{(Conceptual)} Why is fail-closed the correct default for RLS rule misses, and how do you detect them?
    \item \textbf{(Hands-on)} Complete the Thursday project; add an RLS rules dataset restricting a test user to one region and verify exports.
    \item \textbf{(Hands-on)} Build the reverse-ETL loop: \cref{lst:ch16-unload} plus a Lambda pushing to a mock CRM endpoint, with a reconciliation count.
    \item \textbf{(Hands-on)} Inject an anomaly into daily revenue and capture the ML Insights alert end to end.
    \item \textbf{(Design)} Design the identity-resolution strategy for customers who exist in web, app, and store systems with different keys.
    \item \textbf{(Design)} Budget a reverse-ETL sync: 200{,}000 members, API limit 100 calls/s, 100 records/call, business wants 30-minute freshness. Is it feasible? Show the math.
    \item \textbf{(Architecture)} Design the BI layer for a holding company where each subsidiary's data must remain in its own account.
\end{enumerate}

\section{Milestone 4 Capstone: Real-Time Fraud Detection}\label{sec:ch16-capstone}
\index{capstone project}\index{Milestone 4 capstone}\index{fraud detection}

\begin{capstonebox}
\textbf{Mission.} Build an end-to-end real-time fraud-detection platform: simulated card swipes stream into Kinesis, a Managed Flink application flags any card used in two cities within ten minutes (sliding windows, event time), alerts flow through Firehose into Redshift, and a QuickSight dashboard displays live fraud alerts.

\textbf{Estimated time.} 8--10 hours in a sandbox AWS account.

\textbf{What you\textquotesingle{}ll practice.}
\begin{itemize}
    \item Producer design with per-card partition keys (Chapter~13).
    \item Event-time sliding-window detection in Flink SQL (Chapter~15).
    \item Governed delivery with Firehose, raw backup, and idempotent writes (Chapter~14).
    \item A live QuickSight dashboard over the alert stream (this week).
\end{itemize}
\end{capstonebox}

\subsection*{Scenario}
A card network wants same-minute detection of an impossible-travel pattern: one physical card used in two different cities within ten minutes. Swipes arrive as a stream; flagged alerts must reach a fraud-operations dashboard within a minute, with raw events durably archived. The card PAN must never be persisted in plain text anywhere in the analytical estate.

\subsection*{Requirements}
\textbf{Functional requirements.}
\begin{itemize}
    \item A Python producer simulating swipes---including scripted impossible-travel scenarios---with \texttt{card\_id} as the partition key and PANs tokenized at the edge.
    \item A Flink application on event time using a sliding window to detect cross-city use within ten minutes.
    \item Firehose delivery of alerts to Redshift with raw backup to S3; deduplication on \texttt{txn\_id}.
    \item A QuickSight dashboard over the alerts table, refreshing at least every minute.
\end{itemize}

\textbf{Non-functional requirements.}
\begin{itemize}
    \item At least three automated tests: window-logic unit test, late-event handling test, DLQ routing test with a failing case.
    \item Alarms on \texttt{IteratorAgeMilliseconds} and Flink restarts, wired to SNS.
    \item A monthly cost estimate (shard-hours versus MSK, Flink KPUs, Firehose volume) with two documented optimizations.
    \item Security review: KMS-encrypted streams, least-privilege consumers, PAN never persisted in plain text.
\end{itemize}

\subsection*{Reference Architecture}
\begin{figure}[H]
    \centering
    \resizebox{\textwidth}{!}{%
    \begin{tikzpicture}[
        node distance=0.55cm and 1.0cm,
        box/.style={rectangle, rounded corners, draw=violet!70!black, thick, fill=violet!10, align=center, font=\small, minimum height=0.9cm, inner sep=4pt},
        side/.style={rectangle, rounded corners, draw=brown!70!black, thick, fill=brown!10, align=center, font=\footnotesize, inner sep=3pt},
        arr/.style={-{Stealth}, thick}
    ]
    \node[box] (prod) {Producers\\(swipe simulators)};
    \node[box, right=of prod] (kds) {Kinesis Data\\Streams};
    \node[box, right=of kds] (flink) {Managed Flink\\(sliding window)};
    \node[box, right=of flink, yshift=1.0cm] (fh) {Firehose $\rightarrow$ S3\\(raw backup)};
    \node[box, right=of flink, yshift=-1.0cm] (rs) {Redshift\\(fraud alerts)};
    \node[box, right=of rs, yshift=1.0cm] (qs) {QuickSight\\live dashboard};
    \node[side, below=0.6cm of kds] (dlq) {DLQ + idempotent consumers};
    \draw[arr] (prod) -- (kds);
    \draw[arr] (kds) -- (flink);
    \draw[arr] (flink) -- (fh);
    \draw[arr] (flink) -- (rs);
    \draw[arr] (rs) -- (qs);
    \draw[arr, dashed] (dlq) -- (kds);
    \end{tikzpicture}%
    }
    \caption{Milestone 4 reference architecture: real-time fraud detection from swipe simulation to live dashboard.}
    \label{fig:ch16-capstone-arch}
\end{figure}

\subsection*{Implementation Steps}
\begin{enumerate}
    \item Create the stream (on-demand capacity) and the tokenizing producer with scripted fraud scenarios.
    \item Build the Flink SQL application: watermarks on \texttt{event\_ts}, the sliding-window detection query below, alert output to a results stream.
    \item Connect Firehose from the results stream to Redshift with S3 raw backup; dedupe alerts on \texttt{txn\_id}.
    \item Build the QuickSight dashboard on the alerts table with a one-minute refresh and an anomaly insight on alert rate.
    \item Write the three tests, wire the two alarms, run the security review, and build the cost model.
\end{enumerate}

The detection query is the heart of the capstone:

\begin{lstlisting}[language=SQL,caption={Two-cities-in-ten-minutes detection with a hopping window.},label={lst:ch16-fraud-sql}]
SELECT card_id, window_start, window_end,
       COUNT(DISTINCT city) AS city_count
FROM TABLE(HOP(TABLE swipes, DESCRIPTOR(event_ts),
               INTERVAL '10' MINUTES, INTERVAL '1' MINUTES))
GROUP BY card_id, window_start, window_end
HAVING COUNT(DISTINCT city) > 1;
\end{lstlisting}

\subsection*{Deliverables}
\begin{itemize}
    \item Producer, Flink application, Firehose configuration, and dashboard definition, all in version control.
    \item The three automated tests with the failing-case demonstration.
    \item Alarm evidence and a screen recording or screenshots of a scripted fraud event reaching the dashboard.
    \item The cost model with two optimizations and the security review (tokenization verified).
\end{itemize}

\subsection*{Acceptance Criteria}
\begin{itemize}
    \item A scripted cross-city pair flags within the freshness budget; a scripted same-city pair does not.
    \item Late events within the watermark bound attribute correctly; events beyond it route to the side output.
    \item No plain-text PAN exists in any S3 prefix, stream record, or Redshift table---verified by scan.
    \item Alarms fire on injected consumer lag and on a forced Flink restart.
    \item Alert rows in Redshift show no duplicates under producer retries.
\end{itemize}

\subsection*{Stretch Goals}
\begin{itemize}
    \item Add a second detection rule (amount spike versus per-card baseline) as a keyed-process function with per-card state.
    \item Replatform detection on MSK with Kafka consumers and compare operational burden and cost.
    \item Add QuickSight RLS so regional fraud teams see only their region's alerts.
    \item Build the batch reconciliation job that refines streaming alerts with next-day settled data.
\end{itemize}
